% Luc Destombe --- licence CC BY-NC-SA 4.0
% https://creativecommons.org/licenses/by-nc-sa/4.0/deed.fr
% Fichier autonome : compiler avec pdflatex, deux fois (signets, nombre de pages).
\documentclass[a4paper,10pt]{article}
\makeatletter
\newif\iflycee@niveaux
\lycee@niveauxtrue
\newif\iflycee@palatino
\lycee@palatinotrue

\RequirePackage[utf8]{inputenc}
\RequirePackage[T1]{fontenc}
\RequirePackage[french]{babel}
\RequirePackage{etoolbox}

\newcommand{\ShorthandsOff}{\@ifundefined{shorthandoff}{}{\shorthandoff{;:!?}}}
\newcommand{\ShorthandsOn}{\@ifundefined{shorthandon}{}{\shorthandon{;:!?}}}
\ShorthandsOff
\AtBeginDocument{\ShorthandsOn}
\AtBeginEnvironment{tikzpicture}{\ShorthandsOff}

\RequirePackage{geometry}
\RequirePackage{fancyhdr}
\RequirePackage{lastpage}
\RequirePackage{multicol}
\RequirePackage{array}
\RequirePackage{enumitem}
\RequirePackage{xcolor}
\RequirePackage{needspace}

\RequirePackage{amsmath,amssymb}
\RequirePackage{mathpazo}
\DeclareMathSizes{10.95}{10.95}{8}{5.5}
\begingroup\catcode`\_=\active \catcode`\^=\active
\gdef_#1{\sb{\scriptscriptstyle #1}}
\gdef^#1{\sp{\scriptscriptstyle\lycee@moinsepais #1}}
\endgroup
\newcommand\lycee@moins{\mathbin{%
  \setbox\z@\hbox{$\m@th\scriptscriptstyle\lycee@moinsorig$}%
  \dimen@=.55\wd\z@ \advance\dimen@-.5pt
  \hbox to.75\wd\z@{\hss$\m@th\scriptscriptstyle\vcenter{\hbox{%
    \tikz\draw[line width=.5pt,line cap=round](0,0)--(\the\dimen@,0);}}$\hss}}}
\begingroup\lccode`\~=`\- \lowercase{\endgroup\let~}\lycee@moins
\newcommand\lycee@moinsepais{\mathcode`\-="8000 }
\AtBeginDocument{\mathchardef\lycee@moinsorig=\mathcode`\-\relax}
\AtBeginDocument{\catcode`\_=12 \mathcode`\_="8000
  \catcode`\^=12 \mathcode`\^="8000 }
\newcommand\lycee@exposants{%
  \fontdimen13\textfont2=.48\fontdimen6\textfont2
  \fontdimen14\textfont2=.48\fontdimen6\textfont2
  \fontdimen15\textfont2=.48\fontdimen6\textfont2
  \fontdimen13\scriptfont2=.48\fontdimen6\scriptfont2
  \fontdimen14\scriptfont2=.48\fontdimen6\scriptfont2
  \fontdimen15\scriptfont2=.48\fontdimen6\scriptfont2}
\every@math@size\expandafter{\the\every@math@size\lycee@exposants}
\newlength{\retraitcalculs}
\setlength{\retraitcalculs}{2em}

\RequirePackage{tikz}
\usetikzlibrary{arrows.meta,calc,babel,shapes.misc}
\RequirePackage[most]{tcolorbox}
\tcbuselibrary{skins,breakable}

\definecolor{coulDef}{HTML}{1F4E79}
\definecolor{coulProp}{HTML}{7030A0}
\definecolor{coulRem}{HTML}{C55A11}
\definecolor{coulEx}{HTML}{2E7D32}
\definecolor{coulExo}{HTML}{B03A2E}
\definecolor{coulPart}{HTML}{1F4E79}
\definecolor{coulMeth}{HTML}{1B6E4F}
\definecolor{coulCourbe}{HTML}{0B5ED7}
\definecolor{coulCourbeB}{HTML}{C00000}
\definecolor{coulCourbeC}{HTML}{2E7D32}
\definecolor{coulGrille}{HTML}{8C8C8C}
\definecolor{coulRep}{HTML}{1B6E4F}
\definecolor{coulBar}{HTML}{2E7D32}

\newcommand{\R}{\mathbb{R}}
\newcommand{\Rs}{\mathbb{R}^{*}}
\renewcommand{\le}{\leqslant}
\renewcommand{\ge}{\geqslant}
\renewcommand{\approx}{\simeq}
\newcommand{\vect}[1]{\overrightarrow{#1}}
\newcommand{\norme}[1]{\left\lVert #1 \right\rVert}
\newcommand{\intervalle}[2]{\left[#1\,;#2\right]}
\RequirePackage{cancel}
\renewcommand{\CancelColor}{\color{coulExo}}
\newcommand{\fois}[1]{\times{\color{coulExo}#1}}
\newcommand{\barre}[1]{\cancel{#1}}

\tikzset{grille/.style={coulGrille,line width=0.4pt}}
\newcommand{\quadrillage}[4]{\draw[grille,step=1] (#1,#2) grid (#3,#4);}
\tikzset{
  croix/.style={cross out,draw,line width=0.6pt,inner sep=0pt,minimum size=4pt},
  croixdroite/.style={croix,rotate=45,line width=0.8pt},
}
\newcommand{\pt}[3]{\path (#1) node[croix]{} node[#2,font=\small,inner sep=2.5pt]{$#3$};}

\newcommand{\lycee@repere}[4]{%
  \draw[grille,step=1] (#1,#2) grid (#3,#4);
  \draw[-{Stealth[length=2mm]},thick] (#1,0)--({#3+0.4},0) node[below left=-1pt]{$x$};
  \draw[-{Stealth[length=2mm]},thick] (0,#2)--(0,{#4+0.4}) node[below left=-1pt]{$y$};
  \node[below left,font=\scriptsize,inner sep=1.5pt] at (0,0) {$0$};}
\newcommand{\lycee@gradx}[1]{\foreach \k/\l in {#1}{%
  \draw (\k,0.07)--(\k,-0.07) node[below,font=\scriptsize,inner sep=1.5pt]{$\l$};}}
\newcommand{\lycee@grady}[1]{\foreach \k/\l in {#1}{%
  \draw (0.07,\k)--(-0.07,\k) node[left,font=\scriptsize,inner sep=1.5pt]{$\l$};}}
\newcommand{\lycee@intff}[2]{\left[#1\,;#2\right]}
\newcommand{\lycee@intfo}[2]{\left[#1\,;#2\right[}
\newcommand{\lycee@intof}[2]{\left]#1\,;#2\right]}
\newcommand{\lycee@intoo}[2]{\left]#1\,;#2\right[}
\newcommand{\lycee@ens}[1]{\left\{#1\right\}}
\AtBeginDocument{%
  \providecommand{\repere}{\lycee@repere}%
  \providecommand{\gradx}{\lycee@gradx}%
  \providecommand{\grady}{\lycee@grady}%
  \providecommand{\Cf}{\mathcal{C}\sb{\scriptscriptstyle f}}%
  \providecommand{\Cg}{\mathcal{C}\sb{\scriptscriptstyle g}}%
  \providecommand{\intff}{\lycee@intff}%
  \providecommand{\intfo}{\lycee@intfo}%
  \providecommand{\intof}{\lycee@intof}%
  \providecommand{\intoo}{\lycee@intoo}%
  \providecommand{\ens}{\lycee@ens}}

\newlist{questions}{enumerate}{3}
\setlist[questions,1]{label=\arabic*\textdegree),leftmargin=*,itemsep=2pt,topsep=3pt}
\setlist[questions,2]{label=\alph*),leftmargin=*,itemsep=1pt,topsep=2pt}
\setlist[questions,3]{label=\roman*),leftmargin=*,itemsep=1pt,topsep=1pt}
\setlist[itemize]{label=\textbullet,leftmargin=*,itemsep=2pt,topsep=3pt}

\newcommand{\besoinplace}[1]{\par\penalty\@M\begingroup
  \dimen@=#1\relax\dimen@ii=\pagegoal\advance\dimen@ii-\pagetotal
  \ifdim\dimen@>\dimen@ii\ifdim\dimen@ii>\z@\vfil\fi\break\fi\endgroup}

\newcommand{\lycee@pied}{\thepage/\pageref{LastPage}}
\newcommand{\entete}[2]{%
  \pagestyle{fancy}\fancyhf{}%
  \fancyhead[L]{\footnotesize\color{coulPart}\textbf{#1}}%
  \fancyhead[R]{\footnotesize\color{coulPart}\textbf{#2}}%
  \fancyfoot[C]{\footnotesize\lycee@pied}%
  \renewcommand{\headrulewidth}{0.4pt}%
  \renewcommand{\headrule}{\hbox to\headwidth{%
    \color{coulPart!40}\leaders\hrule height \headrulewidth\hfill}}}

\RequirePackage{ccicons}
\newcommand{\lycee@licence}{{\scriptsize\color{gray}%
  \copyright\ Luc Destombe --- \ccbyncsa\ CC BY-NC-SA 4.0}}
\newif\iflycee@licence \lycee@licencetrue
\newcommand{\sanslicence}{\lycee@licencefalse}
\AtBeginDocument{\iflycee@licence\AddToHookNext{shipout/foreground}{%
  \put(0,0){\rlap{\kern\dimexpr1in+\hoffset+\oddsidemargin+\textwidth\relax
    \llap{\raisebox{-\dimexpr1in+\voffset+\topmargin+\headheight+\headsep
      +\textheight+\footskip\relax}{\lycee@licence}}}}}\fi}

\newcommand{\formulesserrees}{%
  \setlength{\abovedisplayskip}{3pt}\setlength{\belowdisplayskip}{3pt}%
  \setlength{\abovedisplayshortskip}{2pt}\setlength{\belowdisplayshortskip}{3pt}}

\setlength{\parindent}{0pt}

\RequirePackage[implicit=false,hidelinks,pdfencoding=auto,psdextra,bookmarksopen,
  bookmarksopenlevel=1,pdfcreator={},pdfproducer={}]{hyperref}
\RequirePackage{bookmark}
\hypersetup{pdfauthor={Luc Destombe},pdfkeywords={CC BY-NC-SA 4.0}}
\newcounter{lycee@signet}
\newcommand{\signet}[3][]{\stepcounter{lycee@signet}%
  \edef\lycee@dest{\if\relax\detokenize{#1}\relax
    signet.\arabic{lycee@signet}\else#1\fi}%
  \hypertarget{\lycee@dest}{}\bookmark[level=#2,dest=\lycee@dest]{#3}}
\pdfstringdefDisableCommands{%
  \def\R{R}\def\vect#1{#1}\def\Cf{Cf}\def\Cg{Cg}\def\fois#1{×#1}%
  \def\barre#1{#1}\def\intervalle#1#2{[#1;#2]}\def\intff#1#2{[#1;#2]}%
  \def\textsuperscript#1{#1}\def\dfrac#1#2{#1/#2}\def\frac#1#2{#1/#2}%
  \def\vec#1{#1⃗}}%
\RequirePackage[official]{eurosym}

\geometry{top=0.8cm,bottom=1.3cm,left=1cm,right=1cm,footskip=0.6cm}

\pagestyle{fancy}
\fancyhf{}
\renewcommand{\headrulewidth}{0pt}
\fancyfoot[C]{\footnotesize Page \thepage{} / \pageref{LastPage}}

\newcommand{\titrefiche}[2]{%
  \begin{center}
    \Large\textbf{#1}\\[2pt]
    \large\textbf{#2}\\[2pt]
  \end{center}
  \vspace{0.10cm}}

\setlength{\columnseprule}{0.4pt}
\setlength{\columnsep}{15pt}
\renewcommand{\columnseprulecolor}{\color{black!45}}
\setlength{\multicolsep}{2pt}

\newcounter{partiefiche}
\renewcommand{\thepartiefiche}{\Roman{partiefiche}}
\newif\ifapresbandeau
\newcommand{\lycee@bandeau}[1]{%
  \par\penalty-600
  \vspace{0.08cm}%
  \noindent\colorbox{coulPart}{%
    \parbox{\dimexpr\linewidth-2\fboxsep\relax}{%
      \color{white}\bfseries\raggedright\strut#1\strut}}%
  \nopagebreak\par\nopagebreak
  \vspace{0.06cm}\nopagebreak
  \apresbandeautrue}
\newcommand{\partiefiche}[1]{%
  \refstepcounter{partiefiche}%
  \lycee@bandeau{\signet{1}{\thepartiefiche{} --- #1}\thepartiefiche{} --- #1}}
\newcommand{\partiefichesansnum}[1]{\lycee@bandeau{\signet{1}{#1}#1}}

\newcounter{exo}
\newlength{\espaceexo}\setlength{\espaceexo}{7pt}
\newlength{\espacetitre}\setlength{\espacetitre}{2pt}
\newcommand{\exo}[1][\relax]{%
  \par
  \ifapresbandeau\apresbandeaufalse\else\penalty-150\addvspace{\espaceexo}\fi
  \refstepcounter{exo}%
  \noindent\signet[exercice-\arabic{exo}]{2}{Exercice \arabic{exo}}%
  \textbf{\color{coulExo}\underline{Exercice \theexo}}%
  \ifx\relax#1\relax\else\textbf{ : #1}\fi
  \par\nopagebreak\vspace{\espacetitre}\nopagebreak
  \ignorespaces}

\setlist[enumerate]{nosep,leftmargin=*,itemsep=2pt,topsep=2pt}
\setlist[itemize]{nosep,leftmargin=*,topsep=2pt}
\newcommand{\choix}[3]{\par\hspace*{1em}%
  \makebox[0.3\linewidth][l]{a) #1}\makebox[0.3\linewidth][l]{b) #2}c) #3}
\newcommand{\deux}[2]{\par\hspace*{0.6em}%
  \makebox[0.48\linewidth][l]{#1}#2}
\newcommand{\trois}[3]{\par\hspace*{0.6em}%
  \makebox[0.32\linewidth][l]{#1}\makebox[0.32\linewidth][l]{#2}#3}

  \renewcommand{\exo}[1][\relax]{%
    \ifapresbandeau\apresbandeaufalse\else\par\penalty-150\fi
    \refstepcounter{exo}%
    \vspace{0.05cm}%
    \noindent\signet[exercice-\arabic{exo}]{2}{Exercice \arabic{exo}}%
    \textbf{\color{coulExo}\underline{Exercice \theexo}}%
    \ifx\relax#1\relax\else\textbf{ : #1}\fi%
    \par\nopagebreak
    \ignorespaces}
  \newcommand{\rappel}[1]{%
    \par\vspace{0.06cm}\nopagebreak
    \noindent\fcolorbox{black!35}{black!5}{%
      \parbox{\dimexpr\linewidth-2\fboxsep-2\fboxrule\relax}{%
        \small\textbf{Rappel.} #1}}%
    \par\vspace{0.06cm}\nopagebreak}
  \setlist[enumerate]{nosep,leftmargin=*}
  \setlist[itemize]{nosep,leftmargin=*}
  \setlength{\multicolsep}{3pt plus 1pt minus 1pt}
  \setlength{\premulticols}{10pt}
  \setlength{\postmulticols}{10pt}
  \newlist{expr}{enumerate}{1}
  \setlist[expr]{label={\Alph*\ $=$},nosep,leftmargin=*,itemsep=0pt}

\renewenvironment{center}{\par\addvspace{3pt}\centering}{\par\addvspace{3pt}}
\formulesserrees
\AtBeginDocument{\formulesserrees}
\clubpenalty=10000
\widowpenalty=10000
\displaywidowpenalty=10000
\brokenpenalty=10000
\setlength{\parskip}{0pt}

\RequirePackage{ragged2e}
\setlength{\RaggedRightRightskip}{0pt plus 1fil}
\AtBeginDocument{\RaggedRight\binoppenalty=10000 \relpenalty=10000 }
\g@addto@macro\@parboxrestore{\RaggedRight}
\makeatother
\geometry{top=0.8cm,bottom=1.3cm,left=1cm,right=1cm,footskip=0.7cm,headheight=12pt}
\usepackage{tkz-tab}

\newtcblisting{python}{%
  listing only,enhanced,
  colback=black!3,colframe=black!35,boxrule=0.5pt,arc=1pt,
  left=4pt,right=4pt,top=1pt,bottom=1pt,
  listing options={language=Python,basicstyle=\ttfamily\small,
    keywordstyle=\color{coulMeth}\bfseries,columns=fullflexible,
    keepspaces=true,showstringspaces=false,aboveskip=0pt,belowskip=0pt}}
\newcommand{\reperesuite}[4]{%
  \draw[grille,step=1] (-1,#2) grid (#1,#3);
  \draw[-{Stealth[length=2mm]},thick] (-1,0)--({#1+0.4},0)
    node[below left=-1pt,font=\footnotesize]{$n$};
  \draw[-{Stealth[length=2mm]},thick] (0,#2)--(0,{#3+0.4})
    node[below left=-1pt,font=\footnotesize]{$#4$};
  \node[below left,font=\scriptsize,inner sep=1.5pt] at (0,0) {$0$};}

\begin{document}
\titrefiche{1\textsuperscript{re} spécialité --- Suites numériques}{Fiche d'exercices du chapitre}

\begin{multicols}{2}

\partiefiche{Définition d'une suite}

\exo[Repérer les termes et les rangs]
Une suite $(u_{n})$, de premier terme $u_{0}$, commence par :
$5$ ;\; $9$ ;\; $13$ ;\; $17$ ;\; $21$ ;\; $25$.
\begin{enumerate}
  \item Donner $u_{0}$ et $u_{3}$.
  \item Quel est le rang du terme $21$ ?
  \item Quel est le 4\textsuperscript{e} terme de la suite ?
  \item Donner le terme qui suit $u_{3}$, puis celui qui le précède.
\end{enumerate}

\exo[Une suite qui commence à $u_{1}$]
Une suite $(v_{n})$, de premier terme $v_{1}$, commence par :
$1$ ;\; $4$ ;\; $9$ ;\; $16$ ;\; $25$.
\begin{enumerate}
  \item Donner $v_{1}$ et $v_{4}$.
  \item Quel est le 5\textsuperscript{e} terme de la suite ? Quel est son rang ?
  \item Conjecturer une expression de $v_{n}$ en fonction de $n$.
\end{enumerate}

\exo[Premier rang]
Pour chaque suite, donner le premier rang $n_{0}$ à partir duquel elle est définie.
\begin{multicols}{2}
\begin{enumerate}[label=\alph*)]
  \item $u_{n}=\dfrac{5}{n}$
  \item $v_{n}=\sqrt{n-3}$
  \item $w_{n}=\dfrac{4}{\sqrt{n-5}}$
  \item $t_{n}=\sqrt{n}+\dfrac{1}{n}$
\end{enumerate}
\end{multicols}

\exo[Les décimales de $\frac{1}{7}$]
On a $\frac{1}{7}=0{,}142857142857\dots$ On note $u_{n}$ le $n$-ième chiffre après la
virgule ; ainsi $u_{1}=1$ et $u_{2}=4$.
\begin{enumerate}
  \item Donner $u_{3}$ à $u_{8}$. Quel est le premier rang de cette suite ?
  \item Calculer $u_{20}$, puis $u_{100}$, sans écrire toutes les décimales.
\end{enumerate}

\partiefiche{Suites définies explicitement}

\exo[Calculer des termes]
Pour chaque suite, calculer le terme de rang $0$, de rang $1$ et de rang $5$.
\begin{multicols}{2}
\begin{enumerate}[label=\alph*)]
  \item $u_{n}=4n-3$
  \item $v_{n}=n^{2}-2n+3$
  \item $w_{n}=\dfrac{3n+2}{n+1}$
  \item $t_{n}=2^{n}-1$
\end{enumerate}
\end{multicols}

\exo[Retrouver un rang]
Pour tout entier naturel $n$, on pose $u_{n}=3n+7$.
\begin{enumerate}
  \item Calculer $u_{10}$.
  \item Pour quel rang $n$ a-t-on $u_{n}=52$ ?
  \item Le nombre $101$ est-il un terme de la suite ? Justifier.
\end{enumerate}

\exo[Exprimer $u_{n+1}$ et $u_{n-1}$]
Pour chaque suite, exprimer le terme de rang $n+1$ et celui de rang $n-1$ en fonction
de $n$ (pour $n\ge 1$).
\begin{multicols}{3}
\begin{enumerate}[label=\alph*)]
  \item $u_{n}=5n-2$
  \item $v_{n}=n^{2}-4n$
  \item $w_{n}=\dfrac{n+3}{2n+1}$
\end{enumerate}
\end{multicols}

\exo[Ne pas confondre]
Pour tout entier naturel $n$, on pose $u_{n}=n^{2}+2n$. Exprimer en fonction de $n$ :
\quad a)~$u_{n+1}$ \quad b)~$u_{n}+1$ \quad c)~$u_{2n}$ \quad d)~$2u_{n}$.

\partiefiche{Suites définies par récurrence}

\exo[Calculer de proche en proche]
Pour chaque suite, calculer les trois termes qui suivent le premier.
\begin{enumerate}[label=\alph*)]
  \item $u_{0}=3$ et, pour tout $n$, $u_{n+1}=2u_{n}-1$.
  \item $v_{0}=2$ et, pour tout $n$, $v_{n+1}=v_{n}+3n-1$.
  \item $w_{1}=16$ et, pour tout $n\ge 1$, $w_{n+1}=\dfrac{w_{n}}{2}+4$.
\end{enumerate}

\exo[Avec un carré, avec une fraction]
Pour chaque suite, calculer les trois termes qui suivent le premier.
\begin{enumerate}[label=\alph*)]
  \item $a_{0}=1$ et, pour tout $n$, $a_{n+1}=a_{n}^{2}+1$.
  \item $b_{0}=1$ et, pour tout $n$, $b_{n+1}=\dfrac{1}{1+b_{n}}$.
\end{enumerate}

\exo[Traduire une situation]
Pour chaque situation, donner le premier terme $u_{0}$ et la relation de récurrence.
\begin{enumerate}[label=\alph*)]
  \item Un salaire de $1\,800$~€ augmente de $2\,\%$ par an ; $u_{n}$ est le salaire
        au bout de $n$ années.
  \item Une ville compte $12\,000$ habitants ; chaque année, $5\,\%$ des habitants
        partent et $900$ personnes arrivent.
  \item Un réservoir contient $500$~L d'eau et en perd $40$~L par jour.
\end{enumerate}

\exo[Une forêt]
Une forêt compte $5\,000$ arbres en 2026. Chaque année, on abat $8\,\%$ des arbres et on
en plante $300$. On note $u_{n}$ le nombre d'arbres en $2026+n$.
\begin{enumerate}
  \item Calculer $u_{1}$ et $u_{2}$.
  \item Justifier que, pour tout entier naturel $n$, $u_{n+1}=0{,}92\,u_{n}+300$.
\end{enumerate}

\partiefiche{Calculatrice}

\exo[Afficher des termes]
Pour tout entier $n$, on pose $u_{n}=n^{2}-10n+30$ ; la suite $(v_{n})$ est définie par
$v_{0}=3$ et $v_{n+1}=2v_{n}-2$.
\begin{enumerate}
  \item Afficher les termes $u_{0}$ à $u_{10}$ et les recopier.
  \item Représenter $(u_{n})$ à la calculatrice avec une fenêtre adaptée.
  \item Donner $v_{10}$ et $v_{20}$.
\end{enumerate}

\exo[Conjecturer une expression]
La suite $(w_{n})$ est définie par $w_{0}=1$ et $w_{n+1}=w_{n}+2n+1$.
\begin{enumerate}
  \item À la calculatrice, afficher $w_{0}$ à $w_{6}$, puis donner $w_{10}$ et $w_{50}$.
  \item Conjecturer une expression de $w_{n}$ en fonction de $n$.
\end{enumerate}

\exo[Représenter et conjecturer]
Pour tout entier $n$, on pose $t_{n}=\dfrac{20}{n+1}$.
\begin{enumerate}
  \item Afficher les termes $t_{0}$ à $t_{9}$ (arrondis à $0{,}01$).
  \item Représenter la suite à la calculatrice et conjecturer son sens de variation.
\end{enumerate}

\exo[Le carré ou la puissance ?]
Pour tout entier $n$, on pose $u_{n}=n^{2}$ et $v_{n}=2^{n}$.
\begin{enumerate}
  \item Afficher côte à côte les termes $u_{0}$ à $u_{10}$ et $v_{0}$ à $v_{10}$.
  \item Pour quelles valeurs de $n$ a-t-on $u_{n}>v_{n}$ ? $u_{n}=v_{n}$ ?
  \item Que peut-on conjecturer pour $n\ge 5$ ?
\end{enumerate}

\partiefiche{Programmer en Python}

\exo[Écrire la fonction]
La suite $(v_{n})$ est définie par $v_{0}=4$ et $v_{n+1}=3v_{n}-5$.
\begin{enumerate}
  \item Écrire une fonction \texttt{terme\_v(n)} qui renvoie $v_{n}$.
  \item Que renvoie \texttt{terme\_v(10)} ?
\end{enumerate}

\exo[Lire un programme]
On considère la fonction ci-dessous.
\begin{python}
def mystere(n):
    u = 5
    for i in range(n):
        u = u/2 + 3
    return u
\end{python}
\begin{enumerate}
  \item Quelle suite cette fonction calcule-t-elle ? Donner son premier terme et sa
        relation de récurrence.
  \item Calculer à la main ce que renvoient \texttt{mystere(1)}, \texttt{mystere(2)} et
        \texttt{mystere(3)}.
\end{enumerate}

\exo[Une relation qui dépend de $n$]
La suite $(w_{n})$ est définie par $w_{0}=0$ et $w_{n+1}=w_{n}+2n+3$.
\begin{python}
def terme_w(n):
    w = 0
    for i in range(n):
        w = ..........
    return w
\end{python}
\begin{enumerate}
  \item Calculer $w_{1}$, $w_{2}$ et $w_{3}$.
  \item Compléter la fonction (dans la boucle, \texttt{i} joue le rôle de $n$).
  \item Que renvoie \texttt{terme\_w(10)} ?
\end{enumerate}

\exo[La suite de Fibonacci]
La suite $(f_{n})$ est définie par $f_{0}=1$, $f_{1}=1$ et, pour tout $n$,
$f_{n+2}=f_{n+1}+f_{n}$ : chaque terme est la somme des deux précédents.
\begin{python}
def fibo(n):
    a = 1
    b = 1
    for i in range(n-1):
        a, b = b, ..........
    return b
\end{python}
\begin{enumerate}
  \item Calculer $f_{2}$ à $f_{6}$.
  \item Compléter la fonction (\texttt{a} et \texttt{b} contiennent deux termes qui se
        suivent) pour qu'elle renvoie $f_{n}$, pour $n\ge 1$.
  \item Que renvoie \texttt{fibo(20)} ?
\end{enumerate}

\partiefiche{Représentation graphique}

\exo[Représenter une suite récurrente]
La suite $(u_{n})$ est définie par $u_{0}=2$ et $u_{n+1}=0{,}5\,u_{n}+2$.
\begin{enumerate}
  \item Dresser le tableau des termes $u_{0}$ à $u_{5}$.
  \item Représenter ces termes dans un repère.
  \item Conjecturer le sens de variation de la suite.
\end{enumerate}

\exo[Représenter une suite explicite]
Pour tout entier $n$, on pose $v_{n}=\dfrac{(n-3)^{2}}{2}$.
\begin{enumerate}
  \item Dresser le tableau des termes $v_{0}$ à $v_{6}$.
  \item Représenter ces termes dans un repère.
  \item Décrire les variations de la suite d'après le graphique.
\end{enumerate}

\exo[Lire un nuage de points]
On a représenté les premiers termes d'une suite $(u_{n})$.
\begin{center}
\begin{tikzpicture}[x=0.45cm,y=0.45cm]
  \reperesuite{7}{0}{5}{u_{n}}
  \gradx{1,2,3,4,5,6} \grady{1,2,3,4}
  \foreach \n/\u in {0/1,1/3,2/4,3/4,4/3,5/1,6/0}
    \path[coulCourbe] (\n,\u) node[croix]{};
\end{tikzpicture}
\end{center}
\begin{enumerate}
  \item Lire $u_{0}$, $u_{2}$ et $u_{5}$.
  \item Pour quels rangs a-t-on $u_{n}=4$ ?
  \item Décrire les variations de la suite sur ces premiers termes.
\end{enumerate}

\exo[Des points sur deux droites]
Pour tout entier $n$, on pose $u_{n}=(-1)^{n}\times n$.
\begin{enumerate}
  \item Calculer $u_{0}$ à $u_{6}$, puis représenter ces termes.
  \item La suite est-elle croissante ? décroissante ?
  \item Sur quelle droite se trouvent les points de rang pair ? ceux de rang impair ?
\end{enumerate}

\partiefiche{Sens de variation : signe de $u_{n+1}-u_{n}$}

\exo[Étudier le signe de la différence]
Étudier le sens de variation de chaque suite.
\begin{enumerate}[label=\alph*)]
  \item $u_{n}=3n^{2}+n$ pour tout entier naturel $n$.
  \item $v_{0}=1$ et, pour tout $n$, $v_{n+1}=v_{n}+n^{2}+2$.
  \item $w_{n}=n^{2}-10n$ pour tout entier naturel $n$.
\end{enumerate}

\exo[Suites décroissantes]
Montrer que chaque suite est décroissante.
\begin{enumerate}[label=\alph*)]
  \item $u_{n}=7-2n^{2}$ pour tout entier naturel $n$.
  \item $v_{0}=5$ et, pour tout $n$, $v_{n+1}=v_{n}-(n-3)^{2}$.
  \item $w_{n}=\dfrac{1}{n+1}$ pour tout entier naturel $n$.
\end{enumerate}

\exo[Conjecturer, puis démontrer]
La suite $(u_{n})$ est définie par $u_{0}=20$ et $u_{n+1}=u_{n}+2n-7$.
\begin{enumerate}
  \item Calculer $u_{1}$ à $u_{5}$ et conjecturer les variations de la suite.
  \item Exprimer $u_{n+1}-u_{n}$ et étudier son signe.
  \item À partir de quel rang la suite est-elle croissante ?
\end{enumerate}

\exo[Un carré toujours positif]
La suite $(v_{n})$ est définie par $v_{0}=0$ et $v_{n+1}=v_{n}+(v_{n}-1)^{2}$.
\begin{enumerate}
  \item Sans calculer de terme, montrer que $(v_{n})$ est croissante.
  \item Calculer $v_{1}$, $v_{2}$ et $v_{3}$. La suite est-elle strictement
        croissante ?
\end{enumerate}

\partiefiche{Sens de variation : quotient $\dfrac{u_{n+1}}{u_{n}}$}

\exo[Comparer le quotient à 1]
Étudier le sens de variation de chaque suite.
\begin{multicols}{3}
\begin{enumerate}[label=\alph*)]
  \item $u_{n}=2\times 5^{n}$
  \item $v_{n}=9\times 0{,}8^{n}$
  \item $w_{n}=\dfrac{4^{n}}{n}$, $n\ge 1$
\end{enumerate}
\end{multicols}

\exo[D'autres quotients]
Étudier le sens de variation de chaque suite.
\begin{multicols}{3}
\begin{enumerate}[label=\alph*)]
  \item $u_{n}=\dfrac{3^{n}}{2^{n}}$
  \item $v_{n}=\dfrac{n}{2^{n}}$, $n\ge 1$
  \item $w_{n}=\dfrac{5^{n}}{n+1}$
\end{enumerate}
\end{multicols}

\exo[Choisir la méthode]
Pour chaque suite, choisir entre la différence et le quotient, puis étudier son sens
de variation.
\begin{multicols}{3}
\begin{enumerate}[label=\alph*)]
  \item $u_{n}=6\times 0{,}5^{n}$
  \item $v_{n}=n^{2}+5n$
  \item $w_{n}=n\times 2^{n}$, $n\ge 1$
\end{enumerate}
\end{multicols}

\exo[Croissante à partir d'un rang]
Pour tout entier $n\ge 1$, on pose $u_{n}=\dfrac{2^{n}}{n^{2}}$.
\begin{enumerate}
  \item Calculer $u_{1}$ à $u_{5}$.
  \item Montrer que $\dfrac{u_{n+1}}{u_{n}}\ge 1$ revient à $n^{2}-2n-1\ge 0$.
  \item En déduire à partir de quel rang la suite est croissante.
\end{enumerate}

\partiefiche{Sens de variation : fonction associée}

\exo[Étudier la fonction associée]
Étudier le sens de variation de chaque suite, définie pour tout entier naturel $n$.
\begin{multicols}{3}
\begin{enumerate}[label=\alph*)]
  \item $u_{n}=n^{3}+6n^{2}+1$
  \item $v_{n}=\dfrac{2n+3}{n+1}$
  \item $w_{n}=n^{3}-3n^{2}$
\end{enumerate}
\end{multicols}

\exo[D'autres fonctions]
Étudier le sens de variation de chaque suite, définie pour tout entier naturel $n$.
\begin{multicols}{3}
\begin{enumerate}[label=\alph*)]
  \item $u_{n}=n^{2}-8n+3$
  \item $v_{n}=5-\dfrac{4}{n+2}$
  \item $w_{n}=\dfrac{n}{n^{2}+1}$
\end{enumerate}
\end{multicols}

\exo[Une erreur fréquente]
La suite $(u_{n})$ est définie par $u_{0}=10$ et $u_{n+1}=f(u_{n})$, avec
$f(x)=0{,}5x+2$. Un élève affirme : « $f$ est croissante, donc $(u_{n})$ est croissante. »
\begin{enumerate}
  \item Calculer $u_{1}$, $u_{2}$ et $u_{3}$.
  \item Que penser de l'affirmation de l'élève ?
\end{enumerate}

\exo[Suite croissante, fonction non croissante]
Pour tout entier $n$, on pose $u_{n}=(n-0{,}4)^{2}$, et $f(x)=(x-0{,}4)^{2}$ sur
$\intfo{0}{+\infty}$.
\begin{enumerate}
  \item Dresser le tableau de variations de $f$ sur $\intfo{0}{+\infty}$.
  \item Calculer $u_{n+1}-u_{n}$ et en déduire le sens de variation de $(u_{n})$.
  \item Une suite croissante a-t-elle forcément une fonction associée croissante ?
\end{enumerate}

\partiefiche{Notion de limite}

\exo[Conjecturer une limite]
Pour chaque suite, calculer les termes indiqués, puis conjecturer sa limite.
\begin{enumerate}[label=\alph*)]
  \item $u_{n}=2+\dfrac{5}{n+1}$ : $u_{9}$, $u_{99}$ et $u_{999}$.
  \item $v_{n}=3n^{2}+1$ : $v_{10}$, $v_{100}$ et $v_{1\,000}$.
  \item $w_{n}=4\times(-1)^{n}$ : $w_{0}$ à $w_{3}$.
\end{enumerate}

\exo[Limite finie ou infinie]
Pour chaque suite, calculer le terme de rang $10$, de rang $100$ et de rang $1\,000$,
puis conjecturer sa limite.
\begin{multicols}{3}
\begin{enumerate}[label=\alph*)]
  \item $a_{n}=5-n^{2}$
  \item $b_{n}=\dfrac{1}{n^{2}}$
  \item $c_{n}=\dfrac{2n+1}{n}$
\end{enumerate}
\end{multicols}

\exo[Avec la calculatrice]
La suite $(u_{n})$ est définie par $u_{0}=1$ et $u_{n+1}=0{,}5\,u_{n}+3$.
\begin{enumerate}
  \item À la calculatrice, donner $u_{10}$, $u_{20}$ et $u_{30}$ (arrondis à
        $10^{-6}$).
  \item Conjecturer la limite de la suite.
\end{enumerate}

\exo[Alterner et converger]
Pour tout entier $n\ge 1$, on pose $u_{n}=\dfrac{(-1)^{n}}{n}$.
\begin{enumerate}
  \item Calculer $u_{10}$, $u_{11}$, $u_{100}$ et $u_{101}$.
  \item La suite change-t-elle de signe ? Semble-t-elle avoir une limite ?
\end{enumerate}

\partiefiche{Problèmes : déterminer un seuil}

\exo[Abonnés d'une newsletter]
Une newsletter compte $3\,000$ abonnés en janvier. Chaque mois, $15\,\%$ des abonnés se
désinscrivent et $800$ nouveaux s'inscrivent. On note $u_{n}$ le nombre d'abonnés au bout
de $n$ mois.
\begin{enumerate}
  \item Calculer $u_{1}$ et $u_{2}$.
  \item Justifier que $u_{n+1}=0{,}85\,u_{n}+800$.
  \item Écrire une fonction \texttt{seuil()} qui renvoie le premier mois où l'on dépasse
        $5\,000$ abonnés.
  \item À l'aide de la calculatrice, donner ce mois.
\end{enumerate}

\exo[Seuil d'une suite explicite]
Pour tout entier $n$, on pose $u_{n}=2n^{2}+3n$.
\begin{enumerate}
  \item Compléter la fonction pour qu'elle renvoie le plus petit $n$ tel que
        $u_{n}>1\,000$.
\begin{python}
def seuil():
    n = 0
    u = 0
    while .......... :
        n = n + 1
        u = ..........
    return n
\end{python}
  \item Retrouver ce rang avec le tableau de la calculatrice.
\end{enumerate}

\exo[Élimination d'un médicament]
Un patient reçoit $40$~mg d'un médicament. Chaque heure, son corps en élimine $20\,\%$.
On note $m_{n}$ la quantité, en mg, présente au bout de $n$ heures.
\begin{enumerate}
  \item Calculer $m_{1}$ et $m_{2}$, puis exprimer $m_{n+1}$ en fonction de $m_{n}$.
  \item Écrire une fonction \texttt{seuil()} qui renvoie le nombre d'heures au bout
        duquel il reste moins de $5$~mg.
  \item Donner ce nombre d'heures.
\end{enumerate}

\exo[Qui dépasse l'autre ?]
Un site compte $500$ abonnés payants et en gagne $40$ par mois ; un autre en compte $200$
et gagne $10\,\%$ d'abonnés par mois. On note $a_{n}=500+40n$ et $b_{n}$ leurs abonnés au
bout de $n$ mois ($b_{0}=200$ et $b_{n+1}=1{,}1\,b_{n}$).
\begin{python}
def depasse():
    b = 200
    n = 0
    while b <= 500 + 40*n :
        b = 1.1*b
        n = n + 1
    return n
\end{python}
\begin{enumerate}
  \item Expliquer ce que renvoie cette fonction.
  \item À l'aide de la calculatrice, donner le résultat et l'interpréter.
\end{enumerate}

\partiefichesansnum{Exercices type bac}

\exo[Automatismes (QCM)]
Pour chaque question, une seule réponse est exacte.
\begin{enumerate}
  \item $u_{n}=n^{2}-1$ ; alors $u_{n+1}$ vaut :\par
        a)~$n^{2}$ \quad b)~$n^{2}+2n$ \quad c)~$n^{2}+1$ \quad d)~$n^{2}+2n+2$
  \item $u_{0}=2$ et $u_{n+1}=3u_{n}-1$ ; alors $u_{2}$ vaut :\par
        a)~$5$ \quad b)~$14$ \quad c)~$17$ \quad d)~$8$
  \item Baisser de $15\,\%$ revient à multiplier par :\par
        a)~$0{,}15$ \quad b)~$1{,}15$ \quad c)~$0{,}85$ \quad d)~$-15$
  \item Si $u_{n+1}-u_{n}=-n^{2}-1$ pour tout $n$, la suite $(u_{n})$ est :\par
        a)~croissante \quad b)~décroissante \quad c)~constante \quad d)~on ne peut pas
        savoir
  \item La limite de $u_{n}=3+\dfrac{1}{n}$ semble être :\par
        a)~$0$ \quad b)~$3$ \quad c)~$4$ \quad d)~$+\infty$
\end{enumerate}

\exo[Une plateforme de covoiturage]
Une plateforme locale de covoiturage compte $600$ conducteurs inscrits en janvier 2026.
Chaque mois, $10\,\%$ des conducteurs se désinscrivent et $40$ nouveaux s'inscrivent. On
note $u_{n}$ le nombre de conducteurs au bout de $n$ mois ; ainsi $u_{0}=600$.
\begin{enumerate}
  \item Calculer $u_{1}$ et $u_{2}$, puis justifier que $u_{n+1}=0{,}9\,u_{n}+40$.
  \item On admet que, pour tout entier $n$, $u_{n}=400+200\times 0{,}9^{n}$. Vérifier
        cette formule pour $n=0$ et $n=1$.
  \item Montrer que $u_{n+1}-u_{n}=-20\times 0{,}9^{n}$. En déduire le sens de
        variation de $(u_{n})$.
  \item Conjecturer la limite de $(u_{n})$. Interpréter.
  \item Compléter la fonction pour qu'elle renvoie le premier mois où il y a moins de
        $450$ conducteurs, puis donner ce mois.
\begin{python}
def seuil():
    u = 600
    n = 0
    while .......... :
        u = ..........
        n = n + 1
    return n
\end{python}
\end{enumerate}

\exo[Coût moyen de fabrication]
Un atelier fabrique $n$ lampes par jour ($n\ge 1$). Le coût moyen de fabrication d'une
lampe, en euros, est $c_{n}=n+\dfrac{400}{n}$.
\begin{enumerate}
  \item Calculer $c_{1}$, $c_{10}$ et $c_{20}$.
  \item On pose $f(x)=x+\dfrac{400}{x}$ sur $\intfo{1}{+\infty}$. Calculer $f'(x)$,
        puis dresser le tableau de variations de $f$.
  \item En déduire le sens de variation de $(c_{n})$ et le nombre de lampes qui rend
        le coût moyen le plus faible.
  \item Pour $n\ge 20$, à partir de quel nombre de lampes le coût moyen dépasse-t-il
        $50$~€ ? On résoudra $n^{2}-50n+400>0$.
\end{enumerate}

\end{multicols}

\end{document}
